\section*{Capítulo \ref{ch:b3:rovibrational-spectroscopy} --- Espectroscopia rotacional e vibracional}

\begin{solution}{exo:b3:rovibrational-spectroscopy:1}
$I = h/(8\pi^2cB) = 6.62607 \times 10^{-34}/(8\pi^2 \times 2.99792 \times 10^{10} \times
1.9225) = \qty{1.4561e-46}{kg.m^2}$. $\mu = \qty{6.85621}{u} = \qty{1.13850e-26}{kg}$, logo
$r_0 = \sqrt{I/\mu} = \qty{113.09}{pm}$.
\end{solution}

\begin{solution}{exo:b3:rovibrational-spectroscopy:2}
$2B_0(J+1)$: 20.88, 41.76, 62.64, \qty{83.52}{cm^{-1}}. $J_{\max} = \sqrt{kT/2hcB} -
\frac12 = \sqrt{208.5/20.88} - 0.5 = 2.7$: o nível $J = 3$ é o mais
povoado, e a raia mais intensa é $4 \leftarrow 3$, em \qty{83.5}{cm^{-1}}.
\end{solution}

\begin{solution}{exo:b3:rovibrational-spectroscopy:3}
A absorção de micro-ondas exige um dipolo permanente: apenas \ce{HCl} e \ce{H2O}.
O Raman rotacional exige uma polarizabilidade anisotrópica: \ce{H2}, \ce{HCl}, \ce{N2},
\ce{CO2}, \ce{H2O}; não os \omterm{def:b3:rovibrational-spectroscopy:rotor-types}{piões esféricos} \ce{CH4} e \ce{SF6}.
\end{solution}

\begin{solution}{exo:b3:rovibrational-spectroscopy:4}
$N_1/N_0 = \eu^{-hc\tilde\nu/kT}$: a \qty{300}{K}, $\eu^{-13.84} = \num{9.8e-7}$; a
\qty{1000}{K}, $\eu^{-4.15} = 0.016$. A \omterm{def:b3:rovibrational-spectroscopy:anharmonic}{banda quente} só se torna perceptível (um por cento
da fundamental) por volta de \qty{1000}{K}.
\end{solution}

\begin{solution}{exo:b3:rovibrational-spectroscopy:5}
$\tilde\omega_e - 2\tilde\omega_ex_e = 2885.3$ e $2\tilde\omega_e - 6\tilde\omega_ex_e =
5665.0$: subtraindo o dobro da primeira da segunda, $-2\tilde\omega_ex_e = -105.6$,
logo $\tilde\omega_ex_e = \qty{52.8}{cm^{-1}}$ e $\tilde\omega_e = \qty{2990.9}{cm^{-1}}$.
\end{solution}

\begin{solution}{exo:b3:rovibrational-spectroscopy:6}
$D_e = 2990.95^2/(4 \times 52.82) = \qty{42341}{cm^{-1}}$; $G(0) = 1495.47 - 13.20 =
\qty{1482.3}{cm^{-1}}$; $D_0 = \qty{40859}{cm^{-1}} = \qty{488.8}{kJ/mol}$, contra o valor verdadeiro
\qty{35760}{cm^{-1}} $= \qty{427.8}{kJ/mol}$: o modelo de Morse superestima em 14\,\%.
\end{solution}

\begin{solution}{exo:b3:rovibrational-spectroscopy:7}
$R(0) - P(2) = 6B_0 = 62.64$: $B_0 = \qty{10.440}{cm^{-1}}$. $R(1) - P(1) = 6B_1 = 60.80$:
$B_1 = \qty{10.133}{cm^{-1}}$. $\alpha_e = \qty{0.307}{cm^{-1}}$. $R(0) = \tilde\nu_0 + 2B_1$ dá
$\tilde\nu_0 = \qty{2885.31}{cm^{-1}}$ (e $P(1) = \tilde\nu_0 - 2B_0$ concorda).
\end{solution}

\begin{solution}{exo:b3:rovibrational-spectroscopy:8}
$\mu_H = \qty{0.97959}{u}$, $\mu_D = \qty{1.90441}{u}$, $\mu_H/\mu_D = 0.51438$. $B_0(\ce{DCl})
= 10.440 \times 0.51438 = \qty{5.370}{cm^{-1}}$. $\tilde\omega_e = 2990.95 \times \sqrt{0.51438}
= 2145.1$, $\tilde\omega_ex_e = 52.82 \times 0.51438 = 27.17$: $\tilde\nu_0 = 2145.1 - 54.3 =
\qty{2090.8}{cm^{-1}}$.
\end{solution}

\begin{solution}{exo:b3:rovibrational-spectroscopy:9}
Deslocamentos $4B_0(J + \frac32)$: 11.94, 19.90, \qty{27.85}{cm^{-1}} (a partir de $J = 0$, 1, 2). Os
dois núcleos \ce{^{14}N} são bósons idênticos de spin 1: os níveis de $J$ par têm seis
estados de spin nuclear, os de $J$ ímpar três, logo as raias que partem de $J$ par
são duas vezes mais intensas que suas vizinhas.
\end{solution}

\begin{solution}{exo:b3:rovibrational-spectroscopy:10}
Um \omterm{def:b3:quantum-model-systems:rotor}{rotor rígido} poria as raias em $\nu_1$, $2\nu_1$, $3\nu_1$: $230542.40$ e
\qty{345813.61}{MHz}, acima das medidas por 4.4 e \qty{17.6}{MHz}. Com
$\nu_{J+1\leftarrow J} = 2B(J+1) - 4D(J+1)^3$: $\nu_1 = 2B - 4D$, $\nu_3 = 6B - 108D$, logo
$3\nu_1 - \nu_3 = 96D$: $D = \qty{0.1835}{MHz}$ e $B = (\nu_1 + 4D)/2 =
\qty{57635.97}{MHz}$. Verificação: $\nu_2 = 4B - 32D = \qty{230538.0}{MHz}$.
\end{solution}

\begin{solution}{exo:b3:rovibrational-spectroscopy:11}
$\Delta G(v + \frac12) = \tilde\omega_e - 2\tilde\omega_ex_e(v+1)$ se anula perto de $v =
\tilde\omega_e/2\tilde\omega_ex_e - 1$; a soma de uma progressão aritmética de $n$ termos
que vai de $\tilde\omega_e - 2\tilde\omega_ex_e$ até quase 0 é cerca de $n\tilde\omega_e/2 \approx
\tilde\omega_e^2/4\tilde\omega_ex_e - \tilde\omega_e/2 \approx D_e - G(0)$. Para \ce{HCl}, os 28 intervalos
($v = 0$ a 27) somam \qty{40857}{cm^{-1}}, contra $D_e - G(0) = \qty{40859}{cm^{-1}}$.
\end{solution}

\begin{solution}{exo:b3:rovibrational-spectroscopy:12}
Só com $J$ par, as transições $J \to J + 2$ partem de $J = 0, 2, 4, \dots$: deslocamentos
$6B, 14B, 22B, \dots$, espaçamento $8B$. \ce{CO2} tem centro de simetria e nenhum
dipolo permanente: não tem espectro de micro-ondas.
\end{solution}

\begin{solution}{pb:b3:rovibrational-spectroscopy:1}
\textbf{1.} $B_0 = \nu/2 = \qty{57.6356}{GHz}$; $B_0/c = \qty{1.92252}{cm^{-1}}$.
\textbf{2.} $\mu = 12 \times 15.99491/27.99491 = \qty{6.85621}{u} = \qty{1.13850e-26}{kg}$.
\textbf{3.} $I = h/8\pi^2B_0 = \qty{1.4560e-46}{kg.m^2}$.
\textbf{4.} $r_0 = \sqrt{I/\mu} = \qty{113.09}{pm}$.
\textbf{5.} $r_0$ excede $r_e = \qty{112.83}{pm}$ em \qty{0.26}{pm}: $B_0$ faz a média de $1/r^2$
sobre a vibração de ponto zero em um poço anarmônico, que passa mais tempo
nas distâncias longas.
\textbf{6.} $2\nu_{1\leftarrow0} = \qty{230.5424}{GHz}$; a raia medida está \qty{4.4}{MHz} abaixo:
distorção centrífuga.
\textbf{7.} $0$, 3.845 e \qty{11.535}{cm^{-1}}, isto é, $E/k = 0$, 5.53 e \qty{16.60}{K}.
\textbf{8.} $N_2/N_1 = \frac53\exp[-(16.60 - 5.53)\,\mathrm K/T]$.
\textbf{9.} $\ln(\frac53/0.85) = 0.673 = 11.06/T$: $T = \qty{16}{K}$.
\textbf{10.} $J_{\max} = \sqrt{16.4/(2 \times 1.43878 \times 1.9225)} - 0.5 = 1.2$: $J = 1$.
\textbf{11.} Um quantum vibracional corresponde a $hc\tilde\nu/k \approx \qty{3080}{K}$: a
\qty{16}{K} nenhuma molécula está vibracionalmente excitada, de modo que a nuvem não
emite luz vibracional.
\textbf{12.} Ela exige apenas \qty{5.5}{K} de excitação, é excitada mesmo no gás mais
frio, e o CO é abundante e tem momento dipolar.
\textbf{13.} $\mu = 13.00335 \times 15.99491/28.99826 = \qty{7.17241}{u}$.
\textbf{14.} $B \propto 1/\mu$: $115.2712 \times 6.85621/7.17241 = \qty{110.1893}{GHz}$.
\textbf{15.} A previsão fica \qty{12}{MHz} (0.011\,\%) abaixo da raia medida: $B_0 = B_e - \frac12\alpha_e$, e o
termo de vibração--rotação $\alpha_e$ varia com outra potência de $\mu$; a
estrutura de equilíbrio é a mesma, a média sobre o ponto zero não.
\textbf{16.} $98.9/1.1 = 90$.
\textbf{17.} A raia de \ce{^{12}C^{16}O} está saturada (opticamente espessa): ela não
cresce mais com a quantidade de gás.
\textbf{18.} A raia rara de \ce{^{13}C^{16}O} permanece proporcional à quantidade de gás, a de
\ce{^{12}C^{16}O} dá a temperatura: juntas, elas medem as duas grandezas.
\textbf{19.} $\tilde\nu_0 = 2169.81 - 2 \times 13.288 = \qty{2143.24}{cm^{-1}}$, \qty{4.666}{\micro m}.
\textbf{20.} $2\tilde\omega_e - 6\tilde\omega_ex_e = \qty{4259.89}{cm^{-1}}$.
\textbf{21.} $G(0) = \frac12\tilde\omega_e - \frac14\tilde\omega_ex_e = \qty{1081.58}{cm^{-1}}$.
\textbf{22.} $R(0) - P(1) = 2(B_0 + B_1) \approx 4B \approx \qty{7.7}{cm^{-1}}$.
\textbf{23.} \ce{^{13}C^{16}O}: \omterm{def:b3:quantum-model-systems:oscillator}{massa reduzida} maior, frequência mais baixa.
\textbf{24.} \textbf{$r_0(\ce{CO}) = \qty{113.1}{pm}$}, a partir de uma única frequência de rádio.
\end{solution}
